\documentclass[11pt]{article}
\usepackage[margin=1in]{geometry}
\usepackage{amsmath}
\usepackage{amssymb}
\usepackage{enumitem}
\usepackage{graphicx}
\usepackage{xcolor}
\usepackage{tcolorbox}

\title{Problem 1a: What Are We Even Doing? \\ (The Confused Student's Guide)}
\author{ECO 310 - Understanding Compensated Demand}
\date{February 23, 2026}

\begin{document}

\maketitle

\section{You're Not Alone - This IS Confusing!}

Let me break this down from scratch. The problem asks for "compensated demand" which sounds scary, but it's actually asking a simple question once you understand what's going on.

\section{Two Different Questions Economists Ask}

There are TWO different optimization problems in microeconomics that look similar but ask opposite questions:

\subsection{Question 1: Regular (Marshallian) Demand}

\begin{tcolorbox}[colback=blue!5!white,colframe=blue!75!black,title=The Normal Question]
\textbf{I have \$100. How do I spend it to be as happy as possible?}

\textbf{Maximize:} Utility $u(x_1, x_2)$ \\
\textbf{Subject to:} Budget constraint $p_1 x_1 + p_2 x_2 \leq m$
\end{tcolorbox}

This is what you do when you go shopping: "I have this much money, let me get the best stuff I can afford."

\subsection{Question 2: Compensated (Hicksian) Demand}

\begin{tcolorbox}[colback=green!5!white,colframe=green!75!black,title=The Backwards Question]
\textbf{I want to be THIS happy (utility level $u^*$). What's the cheapest way to get there?}

\textbf{Minimize:} Expenditure $p_1 x_1 + p_2 x_2$ \\
\textbf{Subject to:} Must reach utility $u(x_1, x_2) \geq u^*$
\end{tcolorbox}

This is like: "I need to be satisfied at level 10. What's the minimum I need to spend?"

\begin{tcolorbox}[colback=red!5!white,colframe=red!75!black,title=KEY INSIGHT]
Problem 1a asks you to solve Question 2 (the backwards one). We're minimizing cost, not maximizing utility!
\end{tcolorbox}

\section{Concrete Example: Coffee vs Tea}

Let me give you a real example to understand this.

\subsection{Setup}

\textbf{Your preferences:} You like coffee and tea equally. One cup of coffee makes you just as happy as one cup of tea. So your utility is:
\begin{equation}
u(\text{coffee}, \text{tea}) = \text{coffee} + \text{tea}
\end{equation}

These are \textbf{perfect substitutes} - you'd trade them 1-for-1 and not care.

\subsection{Current Situation}

\begin{itemize}
\item Coffee costs \$3 per cup ($p_{\text{coffee}} = 3$)
\item Tea costs \$5 per cup ($p_{\text{tea}} = 5$)
\item You want to reach happiness level 10 ($u^* = 10$)
\end{itemize}

\textbf{Question:} What's the cheapest way to get 10 units of happiness?

\subsection{Possible Choices}

You need coffee + tea = 10. Here are some options:

\begin{center}
\begin{tabular}{|c|c|c|c|}
\hline
\textbf{Coffee} & \textbf{Tea} & \textbf{Total Utility} & \textbf{Total Cost} \\
\hline
10 & 0 & 10 & $3(10) + 5(0) = \$30$ \\
8 & 2 & 10 & $3(8) + 5(2) = \$34$ \\
5 & 5 & 10 & $3(5) + 5(5) = \$40$ \\
2 & 8 & 10 & $3(2) + 5(8) = \$46$ \\
0 & 10 & 10 & $3(0) + 5(10) = \$50$ \\
\hline
\end{tabular}
\end{center}

\textbf{Which is cheapest?} Buy 10 coffees, 0 teas for \$30.

\textbf{Why?} Coffee is cheaper! Since you like them equally, always buy the cheaper one.

\begin{tcolorbox}[colback=yellow!10!white,colframe=orange!75!black,title=The Answer for This Example]
Compensated demand: $(x_{\text{coffee}}^c, x_{\text{tea}}^c) = (10, 0)$

This means: "To reach utility 10 cheaply, buy 10 coffees and 0 teas."
\end{tcolorbox}

\section{What If Prices Change?}

\subsection{Example 2: Tea Gets Cheaper}

Now suppose:
\begin{itemize}
\item Coffee costs \$5 per cup ($p_{\text{coffee}} = 5$)
\item Tea costs \$2 per cup ($p_{\text{tea}} = 2$)
\item You still want utility level 10 ($u^* = 10$)
\end{itemize}

Same table:

\begin{center}
\begin{tabular}{|c|c|c|c|}
\hline
\textbf{Coffee} & \textbf{Tea} & \textbf{Total Utility} & \textbf{Total Cost} \\
\hline
10 & 0 & 10 & $5(10) + 2(0) = \$50$ \\
8 & 2 & 10 & $5(8) + 2(2) = \$44$ \\
5 & 5 & 10 & $5(5) + 2(5) = \$35$ \\
2 & 8 & 10 & $5(2) + 2(8) = \$26$ \\
0 & 10 & 10 & $5(0) + 2(10) = \$20$ \\
\hline
\end{tabular}
\end{center}

\textbf{Which is cheapest?} Buy 0 coffees, 10 teas for \$20.

\textbf{Why?} Now tea is cheaper! So buy only tea.

\begin{tcolorbox}[colback=yellow!10!white,colframe=orange!75!black,title=The Answer for Example 2]
Compensated demand: $(x_{\text{coffee}}^c, x_{\text{tea}}^c) = (0, 10)$

This means: "To reach utility 10 cheaply, buy 0 coffees and 10 teas."
\end{tcolorbox}

\section{The General Pattern}

Do you see the pattern?

\begin{itemize}
\item When coffee is cheaper: Buy all coffee, no tea
\item When tea is cheaper: Buy all tea, no coffee
\item When they cost the same: Any mix works (you're indifferent)
\end{itemize}

This makes sense! If you like them equally, always buy the cheaper one.

\section{What Happens When Prices Are Equal? ($p_1 = p_2$)}

This case confuses people, so let's really break it down.

\subsection{Example: Both Goods Cost \$4}

Suppose coffee and tea both cost \$4, and you want utility $u^* = 10$.

You need 10 total units (since $u = x_1 + x_2$). Let's check different combinations:

\begin{center}
\begin{tabular}{|c|c|c|c|}
\hline
Coffee & Tea & Total Utility & Total Cost \\
\hline
10 & 0 & 10 & \$4(10) + \$4(0) = \$40 \\
8 & 2 & 10 & \$4(8) + \$4(2) = \$40 \\
5 & 5 & 10 & \$4(5) + \$4(5) = \$40 \\
3 & 7 & 10 & \$4(3) + \$4(7) = \$40 \\
0 & 10 & 10 & \$4(0) + \$4(10) = \$40 \\
\hline
\end{tabular}
\end{center}

\textbf{Every combination costs exactly \$40!}

\subsection{Why Does This Happen?}

Since both goods:
\begin{itemize}
\item Give 1 utility per unit
\item Cost the same price ($p_1 = p_2 = p$)
\item And you need $u^* = 10$ total units
\end{itemize}

Then any mix of 10 units costs: $p \times 10 = \$4 \times 10 = \$40$

\textbf{Mathematical explanation:}
\begin{align*}
\text{Cost} &= p_1 \cdot x_1 + p_2 \cdot x_2 \\
&= p \cdot x_1 + p \cdot x_2 \quad \text{(since $p_1 = p_2 = p$)} \\
&= p(x_1 + x_2) \\
&= p \cdot u^* \quad \text{(since $x_1 + x_2 = u^*$)}
\end{align*}

No matter how you split $x_1$ and $x_2$, as long as they add to $u^*$, the cost is the same!

\subsection{What Do You Write Down?}

\textbf{For compensated demand:}

When $p_1 = p_2$, the answer is: \textbf{any $(x_1, u^* - x_1)$ where $0 \leq x_1 \leq u^*$}

This means:
\begin{itemize}
\item You can buy $(10, 0)$ (all good 1)
\item You can buy $(7, 3)$ (7 of good 1, 3 of good 2)
\item You can buy $(5, 5)$ (half and half)
\item You can buy $(0, 10)$ (all good 2)
\item Literally \textbf{any split} where the numbers add to $u^* = 10$
\end{itemize}

All are equally optimal!

\textbf{For the expenditure function:}

Even though demand has infinitely many solutions, expenditure is unique:
$$e(p_1, p_2, u^*) = p_1 \cdot u^* = p_2 \cdot u^* = \$4 \times 10 = \$40$$

\begin{tcolorbox}[colback=purple!5!white,colframe=purple!75!black,title=Key Insight]
When $p_1 = p_2$:
\begin{itemize}
\item \textbf{Demand is indeterminate} (infinitely many optimal bundles)
\item \textbf{Expenditure is determinate} (unique cost)
\end{itemize}

This is because every optimal bundle costs the same amount!
\end{tcolorbox}

\section{Writing the Mathematical Answer}

The problem asks you to write this as a general formula. Let's use $x_1$ for good 1 and $x_2$ for good 2 (instead of coffee and tea).

\textbf{Given:}
\begin{itemize}
\item Utility function: $u(x_1, x_2) = x_1 + x_2$
\item Prices: $p_1$ and $p_2$
\item Target utility: $u^*$
\end{itemize}

\textbf{Compensated demand formula:}

\begin{tcolorbox}[colback=yellow!10!white,colframe=orange!75!black,title=Answer to Problem 1a]
\begin{equation}
(x_1^c(p_1, p_2, u^*), x_2^c(p_1, p_2, u^*)) =
\begin{cases}
(u^*, 0) & \text{if } p_1 < p_2 \text{ (good 1 is cheaper)} \\
(0, u^*) & \text{if } p_1 > p_2 \text{ (good 2 is cheaper)} \\
\text{any } (x_1, u^* - x_1) \text{ where } 0 \leq x_1 \leq u^* & \text{if } p_1 = p_2
\end{cases}
\end{equation}
\end{tcolorbox}

\textbf{Translation:}
\begin{itemize}
\item If good 1 is cheaper ($p_1 < p_2$): Buy $u^*$ units of good 1, and 0 of good 2
\item If good 2 is cheaper ($p_1 > p_2$): Buy 0 of good 1, and $u^*$ units of good 2
\item If they cost the same ($p_1 = p_2$): Any combination that adds to $u^*$ works
\end{itemize}

\section{Numerical Check (Using Coffee Example)}

Let's verify our formula works:

\textbf{Example 1:} $p_{\text{coffee}} = 3$, $p_{\text{tea}} = 5$, $u^* = 10$

Since $3 < 5$, coffee is cheaper. Our formula says: $(u^*, 0) = (10, 0)$

Check: 10 coffees + 0 teas. Utility = $10 + 0 = 10$. Cost = $3(10) + 5(0) = \$30$ (cheapest option).

\textbf{Example 2:} $p_{\text{coffee}} = 5$, $p_{\text{tea}} = 2$, $u^* = 10$

Since $5 > 2$, tea is cheaper. Our formula says: $(0, u^*) = (0, 10)$

Check: 0 coffees + 10 teas. Utility = $0 + 10 = 10$. Cost = $5(0) + 2(10) = \$20$ (cheapest option).

\section{Why Is This Called "Compensated" Demand?}

The name is confusing, but here's the idea:

\begin{itemize}
\item In regular demand: "You have \$100, prices change, so you adjust what you buy."
\item In compensated demand: "Prices change, but we imagine giving you just enough money (compensating you) so you can still afford the same happiness level."
\end{itemize}

The "compensation" part means we're holding utility constant at $u^*$, not holding money constant.

\section{Why Do Economists Care About This?}

\textbf{Short answer:} To separate price effects into two parts:

\begin{enumerate}
\item \textbf{Substitution effect:} How do you change consumption when price changes, holding utility constant? (This uses compensated demand!)
\item \textbf{Income effect:} How does the change in your purchasing power affect consumption?
\end{enumerate}

Parts (c)i and (c)ii of the problem ask you to calculate substitution effects, which use your answer to part (a).

\section{Summary: What You Need to Remember}

\begin{tcolorbox}[colback=blue!5!white,colframe=blue!75!black,title=TL;DR]
\textbf{Compensated demand asks:} What's the cheapest way to reach utility level $u^*$?

\textbf{For perfect substitutes} ($u = x_1 + x_2$): Always buy only the cheaper good.

\textbf{Your answer:}
\begin{equation*}
(x_1^c, x_2^c) =
\begin{cases}
(u^*, 0) & \text{if } p_1 < p_2 \\
(0, u^*) & \text{if } p_1 > p_2 \\
\text{any mix adding to } u^* & \text{if } p_1 = p_2
\end{cases}
\end{equation*}

\textbf{Intuition:} If goods are perfect substitutes and you want to minimize cost, buy the cheap one!
\end{tcolorbox}

\section{Connection to Part (b): Expenditure Function}

\subsection{What IS the Expenditure Function?}

The expenditure function answers this question:

\begin{center}
\textit{"What's the MINIMUM amount of money I need to reach utility level $u^*$?"}
\end{center}

It's literally asking: \textbf{How much do I have to spend?}

\subsection{Building the Expenditure Function Step-by-Step}

Let's build this from scratch using our coffee/tea example.

\textbf{Setup:}
\begin{itemize}
\item Coffee costs \$3, tea costs \$5
\item You want utility level 10 (meaning $u^* = 10$)
\item From part (a), we found compensated demand: $(x_1^c, x_2^c) = (10, 0)$
\item Translation: Buy 10 coffees, 0 teas
\end{itemize}

\textbf{Question:} How much does this shopping cart cost?

\textbf{Answer:}
\begin{align*}
\text{Total cost} &= (\text{price of coffee}) \times (\text{coffees bought}) + (\text{price of tea}) \times (\text{teas bought}) \\
&= \$3 \times 10 + \$5 \times 0 \\
&= \$30 + \$0 \\
&= \$30
\end{align*}

That's it! The expenditure function is just the \textbf{total bill} for your optimal shopping cart.

\subsection{Writing It as a Formula}

In general notation:
\begin{equation}
e(p_1, p_2, u^*) = p_1 \cdot x_1^c + p_2 \cdot x_2^c
\end{equation}

Translation:
\begin{itemize}
\item $e$ = expenditure (the total cost)
\item $p_1 \cdot x_1^c$ = (price of good 1) $\times$ (optimal quantity of good 1)
\item $p_2 \cdot x_2^c$ = (price of good 2) $\times$ (optimal quantity of good 2)
\end{itemize}

\textbf{Key insight:} We're using the quantities from compensated demand ($x_1^c, x_2^c$), which we already found in part (a)!

\subsection{Plugging In Our Answer}

From part (a), we know:
\begin{itemize}
\item If $p_1 < p_2$: $(x_1^c, x_2^c) = (u^*, 0)$ (buy only good 1)
\item If $p_1 > p_2$: $(x_1^c, x_2^c) = (0, u^*)$ (buy only good 2)
\end{itemize}

So the expenditure is:

\begin{align}
e(p_1, p_2, u^*) &= p_1 \cdot x_1^c + p_2 \cdot x_2^c \\
&= \begin{cases}
p_1 \cdot u^* + p_2 \cdot 0 = p_1 u^* & \text{if } p_1 < p_2 \\
p_1 \cdot 0 + p_2 \cdot u^* = p_2 u^* & \text{if } p_1 > p_2 \\
p_1 u^* = p_2 u^* & \text{if } p_1 = p_2
\end{cases} \\
&= \min\{p_1, p_2\} \cdot u^*
\end{align}

\subsection{What Does This Mean in English?}

The final formula $e(p_1, p_2, u^*) = \min\{p_1, p_2\} \cdot u^*$ says:

\begin{center}
\textbf{Minimum cost} = (cheaper price) $\times$ (target utility)
\end{center}

\textbf{Why does this make sense?}

Since goods are perfect substitutes:
\begin{enumerate}
\item You need exactly $u^*$ total units to reach utility $u^*$ (because $u = x_1 + x_2$)
\item You'll buy all $u^*$ units of whichever good is cheaper
\item So your cost is: (cheapest unit price) $\times$ (number of units needed)
\end{enumerate}

\subsection{Concrete Examples}

\textbf{Example 1:} Coffee is \$3, tea is \$5, want utility 10
\begin{align*}
e(3, 5, 10) &= \min\{3, 5\} \times 10 \\
&= 3 \times 10 \\
&= \$30
\end{align*}

Why? Coffee is cheaper (\$3 < \$5), so buy 10 coffees at \$3 each = \$30 total.

\textbf{Example 2:} Coffee is \$5, tea is \$2, want utility 10
\begin{align*}
e(5, 2, 10) &= \min\{5, 2\} \times 10 \\
&= 2 \times 10 \\
&= \$20
\end{align*}

Why? Tea is cheaper (\$2 < \$5), so buy 10 teas at \$2 each = \$20 total.

\subsection{The Big Picture}

\begin{tcolorbox}[colback=green!5!white,colframe=green!75!black,title=Understanding the Expenditure Function]
\textbf{What it is:} A function that tells you the minimum cost to reach a given utility level.

\textbf{Inputs:}
\begin{itemize}
\item $p_1, p_2$ = prices of the two goods
\item $u^*$ = target utility level
\end{itemize}

\textbf{Output:} The dollar amount you need to spend

\textbf{How to get it:}
\begin{enumerate}
\item Find compensated demand (part a): What quantities should you buy?
\item Calculate the cost: price $\times$ quantity for each good, then add them up
\end{enumerate}

\textbf{For perfect substitutes:} $e(p_1, p_2, u^*) = \min\{p_1, p_2\} \cdot u^*$
\end{tcolorbox}

\section{Part (c): Understanding Substitution Effects}

\subsection{What Is the Substitution Effect?}

When a price changes, your consumption changes. But WHY does it change? Economists split this into TWO separate reasons:

\begin{enumerate}
\item \textbf{Substitution Effect:} You buy less of the good because it got relatively more expensive (holding utility constant)
\item \textbf{Income Effect:} You buy less because the price increase makes you effectively poorer (can't afford as much)
\end{enumerate}

Part (c) asks ONLY for the substitution effect. This means we ignore the income effect and ask:

\begin{center}
\textit{"If I wanted to stay at my original happiness level, how would my consumption change?"}
\end{center}

\begin{tcolorbox}[colback=purple!5!white,colframe=purple!75!black,title=Key Formula]
\textbf{Substitution Effect} = (New compensated demand) - (Original demand)

Where "new compensated demand" means: the cheapest way to reach the \textbf{original utility} with \textbf{new prices}.
\end{tcolorbox}

This is exactly what you calculated in part (a)! That's why the problem says "using your answer to (a)."

\subsection{Setup from the Problem}

The problem tells us:

\begin{itemize}
\item Initially: $p_1 = 1$, $p_2 = 2$, $m = 10$
\item Consumer buys bundle $(10, 0)$ and earns utility 10
\item Utility function: $u(x_1, x_2) = x_1 + x_2$
\end{itemize}

\textbf{Why $(10, 0)$?} Because good 1 is cheaper (\$1 < \$2), so with perfect substitutes, buy only good 1. With \$10, you can buy 10 units of good 1.

\textbf{Original utility:} $u^* = 10 + 0 = 10$

\subsection{Part (c)(i): Price Increases to 1.5}

\textbf{Question:} Calculate the substitution effect if $p_1$ increases to 1.5.

\subsubsection{Step 1: Set Up}

\begin{itemize}
\item \textbf{Original situation:} $p_1 = 1$, $p_2 = 2$ → bought $(10, 0)$ → utility = 10
\item \textbf{New prices:} $p_1 = 1.5$, $p_2 = 2$
\item \textbf{Target utility:} Keep $u^* = 10$ (the original utility)
\end{itemize}

\subsubsection{Step 2: Find New Compensated Demand}

Use your answer from part (a): what bundle reaches $u^* = 10$ cheaply with new prices?

Check which good is cheaper:
\begin{itemize}
\item Good 1 now costs \$1.5
\item Good 2 still costs \$2
\item Since $1.5 < 2$, good 1 is still cheaper!
\end{itemize}

From part (a), when $p_1 < p_2$:
$$x_1^c(1.5, 2, 10) = (u^*, 0) = (10, 0)$$

So the new compensated demand is $(10, 0)$ - still buy 10 units of good 1, 0 of good 2.

\subsubsection{Step 3: Calculate Substitution Effect}

\begin{align*}
\text{Substitution Effect} &= (\text{New compensated demand}) - (\text{Original demand}) \\
&= (10, 0) - (10, 0) \\
&= (0, 0)
\end{align*}

\begin{tcolorbox}[colback=yellow!10!white,colframe=orange!75!black,title=Answer to Part (c)(i)]
\textbf{Substitution effect:}
\begin{itemize}
\item Change in good 1: $\Delta x_1 = 0$
\item Change in good 2: $\Delta x_2 = 0$
\end{itemize}

The substitution effect is \textbf{zero}!
\end{tcolorbox}

\subsubsection{Why Is It Zero?}

Even though good 1 got more expensive (from \$1 to \$1.50), it's \textbf{still the cheaper good}. Since goods are perfect substitutes, you still buy only good 1 to minimize cost. There's no substitution to good 2 because good 2 is still more expensive.

\textbf{Key insight:} With perfect substitutes, substitution happens only when the price change makes the other good become cheaper!

\subsection{Part (c)(ii): Price Increases to 2.5}

\textbf{Question:} Calculate the substitution effect if $p_1$ increases to 2.5.

\subsubsection{Step 1: Set Up}

\begin{itemize}
\item \textbf{Original situation:} $p_1 = 1$, $p_2 = 2$ → bought $(10, 0)$ → utility = 10
\item \textbf{New prices:} $p_1 = 2.5$, $p_2 = 2$
\item \textbf{Target utility:} Keep $u^* = 10$
\end{itemize}

\subsubsection{Step 2: Find New Compensated Demand}

Check which good is cheaper:
\begin{itemize}
\item Good 1 now costs \$2.5
\item Good 2 still costs \$2
\item Since $2.5 > 2$, good 2 is now cheaper!
\end{itemize}

From part (a), when $p_1 > p_2$:
$$x_1^c(2.5, 2, 10) = (0, u^*) = (0, 10)$$

So the new compensated demand is $(0, 10)$ - buy 0 of good 1, 10 of good 2.

\subsubsection{Step 3: Calculate Substitution Effect}

\begin{align*}
\text{Substitution Effect} &= (\text{New compensated demand}) - (\text{Original demand}) \\
&= (0, 10) - (10, 0) \\
&= (-10, +10)
\end{align*}

\begin{tcolorbox}[colback=yellow!10!white,colframe=orange!75!black,title=Answer to Part (c)(ii)]
\textbf{Substitution effect:}
\begin{itemize}
\item Change in good 1: $\Delta x_1 = -10$ (decrease by 10 units)
\item Change in good 2: $\Delta x_2 = +10$ (increase by 10 units)
\end{itemize}

You completely substitute away from good 1 toward good 2!
\end{tcolorbox}

\subsubsection{Why Such a Large Effect?}

Good 1 was \$1, increased to \$2.5 - that's a huge jump! Now good 1 is MORE expensive than good 2. Since they're perfect substitutes (you like them equally), you switch completely to good 2.

\textbf{Key insight:} With perfect substitutes, the substitution effect is all-or-nothing. You either stay with the same good (if it's still cheaper) or switch completely to the other good (if it becomes cheaper).

\subsection{Summary: Comparing Both Cases}

\begin{center}
\begin{tabular}{|l|c|c|c|}
\hline
\textbf{Scenario} & \textbf{New $p_1$} & \textbf{Which is cheaper?} & \textbf{Substitution Effect} \\
\hline
Original & 1 & Good 1 (\$1 < \$2) & - \\
\hline
Part (c)(i) & 1.5 & Good 1 (\$1.5 < \$2) & $(0, 0)$ - no change \\
\hline
Part (c)(ii) & 2.5 & Good 2 (\$2 < \$2.5) & $(-10, +10)$ - complete switch \\
\hline
\end{tabular}
\end{center}

\begin{tcolorbox}[colback=blue!5!white,colframe=blue!75!black,title=Big Picture]
The substitution effect measures how you respond to relative price changes when trying to maintain your original happiness level.

\textbf{For perfect substitutes:}
\begin{itemize}
\item If the good that got more expensive is \textbf{still cheaper}: No substitution ($\Delta = 0$)
\item If the price increase makes it \textbf{more expensive than the alternative}: Complete substitution (switch entirely to the other good)
\end{itemize}

This is why perfect substitutes have "extreme" substitution effects - it's all or nothing!
\end{tcolorbox}

\end{document}
